\documentclass[10pt,a4paper]{amsart}
\usepackage[margin=19mm]{geometry}
\usepackage{amsmath,amssymb,mathtools}
\usepackage[scr=rsfs,cal=euler]{mathalfa}
\usepackage{xfrac}
\usepackage[unicode,hidelinks,bookmarks=false]{hyperref}
\newcommand{\ie}{i.e.\ }
\newcommand{\cf}{cf.\ }
\newcommand{\resp}{respectively\ }
\newcommand{\Subsection}{\subsection}
\providecommand{\nobreakdash}{\nobreak\hbox{-}}
\setlength{\emergencystretch}{2em}
\multlinegap0pt
\usepackage{amssymb}
\usepackage{mathtools}
\usepackage{tikz}
\newtheorem{definition}{Definition}
\newtheorem{lemma}{Lemma}
\newtheorem{theorem}{Theorem}
\newtheorem{claim}{Claim}
\DeclareMathOperator{\Hilb}{Hilb}
\DeclareMathOperator{\Rays}{Rays}
\DeclareMathOperator{\ZZ}{\mathbf{Z}}
\DeclareMathOperator{\coker}{coker}
\DeclareMathOperator{\dani}{\widetilde{\Omega}}
\DeclareMathOperator{\dual}{\!^{\vee}}
\DeclareMathOperator{\ima}{im}
\DeclareMathOperator{\rank}{rank}
\DeclareMathOperator{\supp}{supp}
\DeclareMathOperator{\tors}{tors}
\title{$K_2$-regularity is Equivalent to Smoothness for Toric Varieties}
\author{Christian Haesemeyer, George Henderson-Walshe}
\date{Source: arXiv:2608.27761v1 (CC BY 4.0)}
\begin{document}
\maketitle

\noindent\emph{Source and licence.} $K_2$-regularity is Equivalent to Smoothness for Toric Varieties. Version arXiv:2608.27761v1 (CC BY 4.0), distributed by arXiv under the Creative Commons Attribution 4.0 International licence, http://creativecommons.org/licenses/by/4.0/, as recorded in the arXiv licence metadata for this exact version and shown on the abstract page https://arxiv.org/abs/2608.27761v1. Full licence notice: \texttt{licenses/arxiv-2608.27761-CC-BY-4.0.txt}.\par\smallskip

\noindent\textit{Editorial scope.} Counted unit: the article's lemma case (source0.tex lines 482--484). The article's complete original proof of it is delivered with it (source0.tex lines 485--573, one \texttt{proof} environment of 89 lines). No mathematics is written, repaired or re-derived: the statement and the proof are the article's own lines, in the article's own notation, and no retained line is substituted or re-flowed.  Retained uncounted as the source-local context the proof needs: lines 388--393; lines 397--399; lines 426--428; lines 475--480.  Nothing was omitted from any retained span, so the package carries no omission record.  No retained line contains a citation, so the wrapper carries no bibliography and no bibliography entry.  No retained line contains a figure environment, an image-inclusion command or drawing code, so no image is delivered and the package declares no asset dependency.  Editorial formatting only, in addition: an AMS \texttt{amsart} preamble that restates the article's own declarations and macros used by the retained lines, namely the environments \texttt{definition}, \texttt{lemma}, \texttt{theorem}, \texttt{claim} on separate counters, together with the standard packages those lines need.  The retained environments print in this document as Definition 1, Lemma 1, Theorem 1, Claim 1 in order of appearance, whereas the article prints the same items with the numbering of its own full document.  The complete native source of the article as distributed in the arXiv e-print for this exact version is bundled unchanged as \texttt{sources/208737/source0.tex}.
\begin{definition} \label{definition: hilbert basis}
For a subset $S$ of a monoid $K$ such that $K \cap (-K) = \{ O \}$, we denote by $\Hilb(S)$ the set of minimal elements with respect to $\leq$, also called the \emph{Hilbert basis} of $S$,
\begin{equation}
\Hilb(S) \coloneq \{x \in S \backslash \{ O \} \mid \forall y \in S, y \leq x \implies y = x \}.
\end{equation}
\end{definition}

\begin{lemma} \label{lemma: hilbert basis face}
If $\sigma$ is a strongly convex cone in the lattice $N$ and $\theta < \sigma\dual$, then $\Hilb(\theta) = \theta \cap H_{\sigma }$. 
\end{lemma}

\begin{lemma} \label{lemma: coker1 = H less R}
	Let $X=U(\sigma,N)$ be an affine toric variety such that $\sigma$ is full-dimensional. If $R_{\sigma} \neq H_{\sigma}$ then $\coker(\beta^1_X) \neq 0$.
\end{lemma}

\begin{theorem} \label{thm:kahler_torsion}
	If $X = U(\sigma,N)$ is a toric variety then for any $m \in S_\sigma$ we have
	\begin{equation}
	\tors (\Omega_{X}^1)(m) = T_m \otimes_{\ZZ} k.
	\end{equation}
\end{theorem}

\begin{lemma} \label{lemma: base case}
    If $X = U(\sigma,N)$ is an affine toric variety such that $\sigma$ is non-simplicial, full-dimensional, and $\rank N = 3$, then either $\ker(\beta^1_X) \neq 0, \coker(\beta^1_X)\neq 0$, or $\coker(\beta^2_X)\neq 0$.
\end{lemma}

\begin{proof}
    We proceed by assuming that $\sigma$ is non-simplicial and that \begin{equation}
        \ker(\beta^1_X)  =\coker(\beta^1_X)= 0
    \end{equation} and deriving that $\coker(\beta^2_X) \neq 0$ by five small combinatorial claims. By Lemma \ref{lemma: coker1 = H less R}, since $\coker(\beta^1_X)=0$, $H_\sigma= R_\sigma$ (as in Definition \ref{definition: hilbert basis}).

    \begin{claim} \label{claim: min support for relation}
        If $\ell \in L(H_\sigma) \backslash \{O \}$, then $|\supp \ell^+|, |\supp \ell^-| \geq 2$.
    \end{claim}
    \begin{proof}
    Since $H_\sigma= R_\sigma$, if $|\supp(\ell^+)|=1$ then one extremal ray is equal to a positive linear combination of the others, which contradicts the definition of an extremal ray.
    \end{proof}

    \begin{claim} \label{claim: dim L + dim Gamma = H}
        If $m \in S_{\sigma}$ then $L((H_\sigma)_m)$ is a free abelian group and $\rank L((H_\sigma)_m) + \dim \Gamma(m)= |(H_\sigma)_m|$.
    \end{claim}
    \begin{proof}
        By definition, $L((H_\sigma)_m)$ is a subgroup of a free abelian group, so it is a free abelian group. For any $n \in \sigma \cap N$, if $\langle m,n\rangle = 0$ then for any $h < m$, since $\langle m - h, n \rangle \geq 0$, we have $\langle h,n \rangle = 0$. Conversely, if $\langle h, n \rangle = 0$ for all $h \in H_\sigma$ then $\langle m, n \rangle = 0$. Therefore $\dim \Gamma(m) = \rank \langle (H_\sigma)_m \rangle = \rank \ima \phi|_{\ZZ^{(H_\sigma)_m}}$. Therefore, by the first isomorphism theorem,
        \begin{align}
            \rank L((H_\sigma)_m) + \dim \Gamma(m) &= \rank \ker \phi|_{\ZZ^{(H_\sigma)_m}} + \rank \phi|_{\ZZ^{(H_\sigma)_m}}\\
            &= \rank \ZZ^{(H_\sigma)_m} = |(H_\sigma)_m|.
        \end{align}
    \end{proof}

    \begin{claim} \label{claim: dim Gamma = 3}
        If $\ell \in L(H_\sigma) \backslash \{ O\}$ then $\dim \Gamma(\psi(\ell)) = 3$.
    \end{claim}
    \begin{proof}
        Let $m = \psi(\ell)$. If $m \in \theta$ for some dual face $\theta < \sigma\dual$ then by Lemma \ref{lemma: hilbert basis face}, $(H_\sigma)_m \subseteq \Hilb(\theta) =  \theta \cap H_\sigma$. Since $H_\sigma = R_\sigma$, we have $(H_\sigma)_m \subseteq \theta \cap R_\sigma= \Rays(\theta)$. However, if $\dim \theta \leq 2$, then $\theta$ must be simplicial, so $|\Rays(\theta)| = \dim \theta$. Therefore $|(H_\sigma)_m| \leq \dim \theta \leq \dim \Gamma(m)$. Therefore by Claim \ref{claim: dim L + dim Gamma = H}, $\rank L((H_\sigma)_m) = 0$, which contradicts the assumption that $\ell \in L(H_\sigma)\backslash \{O \} $ and $m = \psi(\ell)$. Thus if $\ell \in L(H_\sigma) \backslash \{ O\}$ then $\dim \Gamma(\psi(\ell)) \geq 3$. Since $\rank N =3$, a cone in $M$ may be of dimension at most $3$, so $\dim \Gamma(\psi(\ell)) = 3$.  \end{proof}

    \begin{claim} \label{claim: dim L(Hm) = 1}
        For any $m \in \Hilb(\psi(L(H_\sigma)))$, $\rank L((H_\sigma)_m) = 1$.
    \end{claim}
    \begin{proof}
        Since $\ker(\beta^1_X)=0$, Theorem \ref{thm:kahler_torsion} implies that $T_m = 0$ for all $m \in S_\sigma$. Therefore
	
	\begin{equation}
		L((H_\sigma)_m) = \langle \ell \in L((H_\sigma)_m) \mid \psi(\ell) \leq m \rangle,
	\end{equation}
	
	If, moreover, $m \in \Hilb(\psi(L(H_\sigma)))$ then $\psi(\ell) \leq m$ implies $\ell \in \psi^{-1}(m)$, so
	
	\begin{equation}
		L((H_\sigma)_m) = \langle \ell \in L((H_\sigma)_m) \mid \psi(\ell) \leq m \rangle = \langle \psi^{-1}(m) \cap \ZZ^{(H_\sigma)_m}\rangle.
	\end{equation}

    Let $s$ denote $\rank \langle \phi^{-1}(m) \cap \ZZ^{(H_\sigma)_m} \rangle - 1$, and let $\{ b_0, \ldots, b_s \}$ be a basis for the sublattice $ \langle \phi^{-1}(m) \cap \ZZ^{(H_\sigma)_m} \rangle$. Then we have 
    \begin{equation}
        \langle \psi^{-1}(m) \cap \ZZ^{(H_\sigma)_m}\rangle = \langle b_i - b_j \mid 0 \leq i,j\leq s \rangle = \langle b_0 -b_i \mid 1 \leq i \leq s \rangle.
    \end{equation}
    
    Since the $b_i$ are independent for $1 \leq i \leq s$, the $(b_0-b_i)$ are also independent, so $\rank L((H_\sigma)_m) = \rank \langle b_0 - b_i \mid 1 \leq i \leq s \rangle = s$. Therefore by Claims \ref{claim: dim L + dim Gamma = H} and \ref{claim: dim Gamma = 3},
    \begin{equation} \label{eq: s + 3 = H}
        s+3=s + \dim \Gamma(m) = |(H_\sigma)_m|.
    \end{equation}

    Next, if there exists $i \neq j$ and $h \in \supp(b_i) \cap \supp(b_j)$ then $e_h - b_i ,e_h - b_j  \in \ZZ^{(H_\sigma)_m+}$, so $\psi(e_h - b_i -(e_h - b_j ))=m-h < m$, which contradicts $m \in \Hilb(\psi(L(H_\sigma)))$. Therefore $\supp(b_i) \cap \supp(b_j) = \emptyset$ for all $i \neq j$, so
    \begin{equation}
        |(H_\sigma)_m| = \sum^s_{i=0} |\supp b_i|.
    \end{equation}
    By Claim \ref{claim: min support for relation}, for each $i \neq j$, $|\supp(b_i - b_j)| \geq 4$. Therefore we have
    \begin{align}
        2|(H_\sigma)_m| &= 2 \sum_{i=0}^s|\supp b_i| \\
        &= |\supp(b_{s} - b_0)| + \sum_{i=0}^{s-1} | \supp (b_i - b_{i+1})|\\
        & \geq 4(s+1).
    \end{align}

    Since $s+3 = |(H_\sigma)_m|$, this implies $2(s+3) \geq 4(s+1)$, so $s \leq 1$. But $s \geq 1$, so $\rank L((H_\sigma)_m)=s=1$.
    \end{proof}

    \begin{claim} \label{claim: coker 2 non zero}
        If $m \in S_\sigma$, $\dim \Gamma(m) = 3$ and $\rank L((H_\sigma)_m)=1$, then 
        \begin{equation}
            \coker(\beta_X^2)(m) \neq 0
        \end{equation}
    \end{claim}
    \begin{proof}
        Since $\rank L((H_\sigma)_m)=1$, there is a unique (up to sign) generator of $L((H_\sigma)_m)$, $\ell$. By Claim \ref{claim: dim L + dim Gamma = H}, $|(H_\sigma)_m|= \dim \Gamma(m) + \rank L((H_\sigma)_m) = 3 + 1 = 4$. By Claim \ref{claim: min support for relation}, $|\supp\ell^+|, |\supp \ell^-| \geq 2$, so
        \begin{equation}
            2 + 2 \leq |\supp\ell^+|+ |\supp\ell^-| = |\supp\ell|  = |(H_\sigma)_m| = 4.
        \end{equation}
        So we must have $|\supp\ell^+| = |\supp \ell^-| = 2$. Let $e_1, e_2 = \supp(\ell^+)$ and $\{e_3,e_4 \} = \phi( \supp \ell^-)$. Since $\phi(e_i) \in H_\sigma$ for each $1 \leq i \leq 4$, and the $e_i$ are distinct, the $\phi(e_i)$ are also distinct. Therefore we have
        \begin{equation}
            \Omega^2_X(m) = \langle \chi^{m-\phi(e_1)-\phi(e_2)} d\chi^{\phi(e_1)}\wedge d\chi^{\phi(e_1)}, \chi^{m-\phi(e_3)-\phi(e_4)} d\chi^{\phi(e_3)}\wedge d\chi^{\phi(e_4)} \rangle,
        \end{equation}
        so $\dim \Omega^2_X(m) = 2$. On the other hand, $\dim \dani^2_X(m) = \dim \Gamma(m)=3$, so $\dim \Omega^2_X(m) < \dim \dani_X^2(m)$. Since $\beta^2_X$ is an $M$-graded linear map, $\coker(\beta^2_X)(m) \neq 0$. 
    \end{proof}

We can now complete the proof of the Lemma using claims \ref{claim: dim Gamma = 3}, \ref{claim: dim L(Hm) = 1}, and \ref{claim: coker 2 non zero}. Since $\sigma$ is not simplicial, $L(H_\sigma) \neq 0$, so there exists $m \in \Hilb(\psi(L(H_\sigma)))$. By Claim \ref{claim: dim Gamma = 3}, $\dim \Gamma(m)=3$ and by Claim \ref{claim: dim L(Hm) = 1}, $\rank L((H_\sigma)_m)=1$. Therefore the hypotheses of Claim \ref{claim: coker 2 non zero} are satisfied, so $\coker(\beta_X^2)\neq 0$.
\end{proof}

\end{document}
